\section{Averaging and propagation on the hidden complement}
\label{sec:hidden}

Assume $0<\beta<1$. Our objective is to bound propagation on the
complement of the recipe space. This bound separates the bulk spectral
edge from any slower eigenvalues visible in that space.
We use the typicality convention of \Cref{thm:path-calculus} for the
path law, endpoint vectors, and auxiliary trees throughout. Thus the
initial density is bounded above by $\exp(o(n))$ relative to equilibrium,
the pointwise bounds are polynomial in $n$, and dimension tends to
infinity before any approximation or large-time limit.
An endpoint vector $u$ is \emph{hidden} if
\[
 \E\abs{u^t B}^2\longrightarrow0
\]
for every fixed recipe vector $B$ at that endpoint. We write
$\norm{u}_2=(\E\norm{u}^2)^{1/2}$; the norm inside the expectation is
the unnormalized Euclidean norm. The same notation for a diagonal label
means $\norm{a}_{2,\tau}=\tau(a^2)^{1/2}$, and will always carry the
subscript $\tau$ to distinguish these two normalizations.

\begin{theorem}[Hidden propagation]\label{thm:hidden-propagation}
Let $\lambda_{\rm sp}$ be the spin spectral edge in
\Cref{prop:spin-spectrum}. For every $z<\lambda_{\rm sp}$ there is a
finite constant $C_z$ with the following property. Fix $t\geq1$, and let
$u=u(X_0)$ and $p=p(X_t)$ be polynomially bounded endpoint vectors with
bounded second moments. Suppose that $u$ is hidden, $p=df$ is an exact
gradient, and
\begin{equation}\label{eq:hidden-energy-hypothesis}
 \E\sum_j w_j(X_t)\norm{d_jp(X_t)}^2\leq C_t<\infty
\end{equation}
uniformly along the dimension sequence. Then
\begin{equation}\label{eq:hidden-propagation}
 \limsup_{n\to\infty}\abs{\E u^tK_t p}
 \leq C_z e^{-zt}
       \limsup_{n\to\infty}\bigl(\norm{u}_2\norm{p}_2\bigr).
\end{equation}
The constant does not depend on the bound in
\Cref{eq:hidden-energy-hypothesis}; that bound is allowed to depend on
the fixed time $t$. Equivalently, the rate on the right is
$\exp[-(\lambda_{\rm sp}-o(1))t]$, with the dimension limit taken first.
The assertion also holds in the integrated deterministic-time and mask
comparisons of \Cref{thm:path-calculus,thm:row-projected-moments}.
\end{theorem}

We first average a Wick string from its two endpoints, with at most
one vertex retaining its actual path label. We then control the endpoint
predictors by likelihood caps and sum the positive rooted excursions
at every parameter strictly below the spin edge. The derivative work
belongs to the averaging step: exactness converts
\Cref{eq:hidden-energy-hypothesis} into an unweighted Hessian bound
after a field cutoff.

\subsection{Wick strings and two-ended centering}
\label{subsec:hidden-wick}

Begin with fixed times, finitely many hops, and bounded diagonal labels.
In the word calculus, put $s=J/\beta$ and denote by
$B_k(a_0,\ldots,a_k)$ the Wick string with $k$ unpaired $s$-edges
and vertex labels $a_0,\ldots,a_k$. The normalization is
$B_0(a)=D_a$, and contracting adjacent edges around a vertex labeled
$a$ replaces that label by $\tau(a)$. Powers of $\beta$ are kept outside
this notation, so each actual path edge contributes a factor $\beta$.
Removing complete noncrossing excursions expresses a general path word
as a sum of Wick strings. At each surviving vertex, multiply the labels
from its repeated visits; keep the removed excursions as separate scalar
trace factors. The surviving labels have ordered time spans with
disjoint interiors.

In the Fock representation from \Cref{thm:path-calculus},
$s=l+l^*$ and $l^*D_al=\tau(a)I$. The Wick string is the sum of the
$k+1$ terms having creations to the left of a cut and annihilations to
its right. In particular,
\begin{equation}\label{eq:hidden-cut-bound}
 \norm{B_k(a_0,\ldots,a_k)}_{\op}
 \leq\sum_{h=0}^k \norm{a_h}_{\infty}
                     \prod_{j\ne h}\norm{a_j}_{2,\tau}.
\end{equation}
The norm here is taken in the representation with the full diagonal
base. A diagonal with support of limiting trace zero can still act on
that base. If it occurs at the cut, \Cref{eq:hidden-cut-bound} retains
its supremum norm.

There is also the exact splitting identity, for $0<r<k$,
\begin{align}
 B_k(a_0,\ldots,a_k)
 ={}&B_r(a_0,\ldots,a_{r-1},1)D_{a_r}
          B_{k-r}(1,a_{r+1},\ldots,a_k)\notag\\
 &-\tau(a_r)B_{r-1}(a_0,\ldots,a_{r-1})
          B_{k-r-1}(a_{r+1},\ldots,a_k).
 \label{eq:hidden-wick-split}
\end{align}
No subtraction is needed at an end vertex. For an interior vertex,
expand the Wick polynomial by inclusion--exclusion over disjoint
adjacent contracted pairs. The first line collects the choices that
leave the pair incident to $r$ uncontracted. The remaining choices
contract that pair and carry a minus sign, producing the second line.
This verification allows arbitrary vertex labels; constancy is not used.

For a label $a$ on a path from $0$ to $t$, let $a^0$ and $a^1$ denote
its stationary conditional predictors from $X_0$ and $X_t$,
respectively. They are the tracial predictors supplied by
\Cref{thm:path-calculus}, represented by bounded approximating recipes
in successive limits. Covariance identities require a different order:
choose a bounded path approximant and subtract its exact conditional
prediction. Keeping these operations distinct avoids a pointwise
regularity requirement on an exact quenched predictor at rare states.

\begin{lemma}[Two-ended averaging]\label{lem:hidden-averaging}
Let $a_0,\ldots,a_k$ be bounded path labels of the preceding kind.
Under the endpoint hypotheses of \Cref{thm:hidden-propagation},
\begin{align}
 B_k(a_0,\ldots,a_k)\ \sim\ 
 &\sum_{r=0}^k
 B_k(a^0_0,\ldots,a^0_{r-1},a_r,a^1_{r+1},\ldots,a^1_k)\notag\\
 &-\sum_{r=0}^{k-1}
 B_k(a^0_0,\ldots,a^0_r,a^1_{r+1},\ldots,a^1_k).
 \label{eq:hidden-averaging}
\end{align}
Here $T\sim T'$ means that $\E u(X_0)^t(T-T')p(X_t)$ tends to zero
after the prescribed successive approximations. The identity is an
identity in these bilinear tests, not an operator-norm replacement of
every path label by its predictor. The error will contain two
distinguished factors: one centered from the initial endpoint and one
from the terminal endpoint.
\end{lemma}

To prove the averaging identity, first isolate its algebraic error.
Multilinearity gives the exact difference between the two sides of
\Cref{eq:hidden-averaging}:
\begin{equation}\label{eq:hidden-double-centered-error}
 \sum_{0\le r<q\le k}
 B_k\bigl(a^0_{<r},\ a_r-a_r^0,\ a_{r<j<q},\
                    a_q-a_q^1,\ a^1_{>q}\bigr).
\end{equation}
Telescope at the first mismatch from the left and then telescope the
remaining suffix from the right. The all-initial term cancels the
last term of the second telescope. Thus
\Cref{eq:hidden-double-centered-error} retains the actual labels
between the two distinguished vertices. It remains to prove that each
doubly centered summand has zero limiting bilinear test.

The first use of the next lemma is to cut off large terminal fields
before estimating the Hessian. A localized endpoint vector can give
those fields substantial Euclidean mass even when their spatial trace
is small. The two centered factors make their removal possible in the
present bilinear test.

\begin{lemma}[Small-trace replacements in a centered string]
\label{lem:hidden-small-trace}
Consider a summand of \Cref{eq:hidden-double-centered-error}. Its
bounded labels may be approximated in tracial $L^2$ while keeping the
two distinguished factors centered at their respective endpoints.
The error in the hidden bilinear test tends to zero with the tracial
errors. A bounded diagonal of vanishing trace support can also be
inserted at the terminal endpoint at zero cost in this test.
\end{lemma}

\begin{proof}
The Fock comparisons in this proof first remove the terminal visible
component, as in the endpoint projection argument used below for the
predictor caps.  Fix a finite family of endpoint recipes with nonsingular
limiting Gram matrix, and let $Q_m$ be their actual orthogonal projection
on the corresponding Gram-good event, and zero off that event.
Both $Q_m$ and $I-Q_m$ are contractions.  For every fixed capped path word
$W$, Cauchy--Schwarz bounds the removed pairing by
$\|Q_mW^tu\|_2\|p\|_2$, which tends to zero at fixed $m$ by
Assertion~\textup{(ii)} of \Cref{thm:path-calculus} and the fixed inverse
Gram bound.  The weighted typicality convention removes the exceptional
Gram event.  Exhaust the endpoint recipe completion, omitting null
directions, and choose $m=m(n)\to\infty$ sufficiently slowly that these
errors vanish for every fixed word and $(I-Q_m)p$ is hidden for every
fixed endpoint recipe.  Its second-moment and polynomial bounds follow
from contraction.  The projection depends only on the shared endpoint
and its fixed seeds, so permutations of the replicated branches preserve
it.  This projection is used only for the Fock bilinear comparisons;
we retain the original exact $p$ for the Hessian estimates that follow.

Work on one finite tree and use the cut expansion behind
\Cref{eq:hidden-cut-bound}. Every cut away from the replaced factor
puts that factor on a pure creation or pure annihilation side. Its trace
two-norm then bounds the error, multiplied by the fixed bounds on the
other factors and the two endpoint norms. Only the cut at the replaced
factor needs another argument. At that cut, at least one distinguished
centered factor lies on the pure side leading to its own centering
endpoint.

Attach $b$ conditionally independent copies at this endpoint. Lift
this single replicated star to the Fock representation, and symmetrize
by taking a direct sum over permutations of the copies. Repeatedly
applying $l^*D_al=\tau(a)I$ expresses inner products of pure-side strings
as products of pair traces. For different copies, one factor is the
pair trace of the centered label and has zero limit. With approximated
predictors, bound its absolute value by $\eta$, tending to zero with
the approximation. Write $C_\alpha$ for the pure-side creation operator
on copy $\alpha$, oriented toward the shared endpoint. The Gram bounds
are
\[
 \norm{C_\alpha^*C_\alpha}\leq C,\qquad
 \norm{C_\alpha^*C_\gamma}\leq C\eta\quad(\alpha\ne\gamma),\qquad
 \norm{[C_1\ \cdots\ C_b]}^2\leq C(1+(b-1)\eta).
\]
If $v$ denotes the vector at the shared endpoint, symmetry of the
lifted star gives
\[
 \E\norm{C_1^*v}^2
 =\frac1b\sum_{\alpha=1}^b\E\norm{C_\alpha^*v}^2
 \leq C(b^{-1}+\eta)\E\norm{v}^2.
\]
The vector at the opposite end may depend on its whole branch;
Cauchy--Schwarz now bounds the squared bilinear contraction by
\[
 C\bigl(b^{-1}+\eta\bigr)\norm{u}_2^2\norm{p}_2^2.
\]
The fixed caps bound the diagonal Gram blocks. Apply Cauchy--Schwarz
and take the limits $n\to\infty$, $\eta\to0$, and $b\to\infty$ in
that order to obtain zero. All cuts for this replacement are treated
on the same star. The argument makes no simultaneous choice of
creation operators on stars rooted at different endpoints.

A terminal insertion with small trace support has the same two cases.
Every nonterminal cut pays its trace two-norm; at the terminal cut,
the left distinguished factor lies on its centered pure side, so the
replication argument applies. Iterate for finitely many replacements.
When the replaced factor is an approximated distinguished label,
keep its exact forward or backward subtraction throughout this step
and use the other distinguished factor at the exceptional cut.

The word comparison retains every rare diagonal's action on the base.
Small trace is used on a pure side, and centering treats the remaining
cut; neither step uses the false implication
$\tau(c^2)\to0\Rightarrow\norm{D_c}_{\op}\to0$.
\end{proof}

\subsection{The Hessian cutoff and the covariance norm}
\label{subsec:hidden-covariance}

The covariance proof will pair a terminal Hessian with a matrix in
Hilbert--Schmidt norm. We first obtain the required unweighted Hessian
bound. Let $h_i(x)=(Jx)_i$ and choose a smooth function $\chi_M$ equal to one
on $[-M,M]$, vanishing outside $[-M-1,M+1]$, and with uniformly bounded
first derivative. In each double-centered test, replace $p$ by
\[
 p^{(M)}_i(x)=\chi_M(h_i(x))p_i(x).
\]
As $M\to\infty$, the field tail has vanishing spatial trace.
\Cref{lem:hidden-small-trace} therefore justifies this replacement
inside each double-centered test. This use of $p^{(M)}$ does not
require unweighted-norm approximation of an arbitrary localized $p$.

Write $dp=(d_jp_i)_{i,j}$, with the component index as row and the
derivative index as column. The exact product rule is
\begin{equation}\label{eq:hidden-product-rule}
 d_j(fg)=f\,d_jg+g\,d_jf-2x_j(d_jf)(d_jg),
 \qquad f(x^j)=f(x)-2x_jd_jf(x).
\end{equation}
Using the second form of this rule gives
\[
 d_jp_i^{(M)}
 =\chi_M(h_i(x^j))d_jp_i+p_i\,d_j\chi_M(h_i).
\]
For sufficiently large $n$, $\max_{ij}|J_{ij}|\leq1$. On the support
of the first coefficient, $|h_i(x)|\leq M+3$, and hence
$w_i(x)\geq c_M>0$. Mixed-gradient symmetry and the row square bound
therefore give the pointwise estimate
\begin{align}
 \norm{dp^{(M)}}_{\HS}^2
 &\leq C_M\sum_{i,j}w_i|d_ip_j|^2
       +C\sum_i p_i^2\sum_jJ_{ij}^2\notag\\
 &\leq C_M\sum_i w_i\norm{d_ip}^2+C\norm{p}^2.
 \label{eq:hidden-cutoff-hessian}
\end{align}
The cutoff supplies the bounded unweighted mean-square Hessian
needed below. This conversion is the averaging proof's only use of
the symmetry $d_jp_i=d_ip_j$.

We will absorb fixed smooth endpoint words into the terminal vector,
so we also need preservation of this Hessian bound. Check the two
generators, a bounded diagonal $D_f$ and $J$. For a diagonal, write
$Df=(d_jf_i)_{i,j}$ and use
\begin{equation}\label{eq:hidden-matrix-product-rule}
 d(D_fq)=D_f\,dq+D_q\,Df-2(Df\circ dq)D_x.
\end{equation}
The derivative matrix $Df$ has bounded operator norm by the smooth
recipe calculus. Consequently
$\norm{D_qDf}_{\HS}\leq\norm{q}\norm{Df}_{\op}$, and the other two
terms are controlled by $\norm{dq}_{\HS}$. Multiplication by $J$ costs
its bounded operator norm. These estimates, with the exact last term
of \Cref{eq:hidden-matrix-product-rule}, prove the assertion by
induction through the endpoint word.

The second preparation concerns a correlated hidden vector. We need a
pointwise operator estimate for a conditionally averaged derivative,
because a tracial or averaged Hilbert--Schmidt error need not remain
small after its rows are weighted by that vector.

At a joining time $s$, condition on $X_s=x$ and the realized past $\pi$.
For a fixed smooth multitime site-vector calculation $F$ satisfying
\Cref{lem:path-mixed-tensors}, put
\[
 A(x,\pi)=d_x\E[F\mid X_s=x,\pi],
\]
with output as row and derivative index as column. First take the
conditional expectation. \Cref{lem:path-conditional-derivative}
then supplies common linear truncations $A_N$ with
\begin{equation}\label{eq:hidden-conditional-op}
 \norm{(A-A_N)^tc}(x,\pi)
 \leq C_F\varepsilon_N\norm c,\qquad
 \sup_{x,\pi}\norm{A-A_N}_{\op}\leq C_F\varepsilon_N,
 \quad\varepsilon_N\downarrow0.
\end{equation}
The first inequality holds for every output vector $c$ at the same
fixed $(x,\pi)$, using the same truncation. Taking its supremum over
$c$ therefore gives the operator estimate before any expectation.

Recall why this lemma applies to multitime labels. Their mixed tensors
retain at most one difference in each observed-state slot. When a
conditional starting state $z$ is copied into a stored parent $y$, the
exact rule is
\[
 d_j[G(y,z)|_{z=y}]
 =[d_j^yG+d_j^zG-2y_jd_j^yd_j^zG]_{z=y}.
\]
The mixed correction is among those retained tensors. The finite
leaf-elimination argument of \Cref{lem:path-conditional-derivative}
uses point-start $K$ bounds on the evolving index, with all other
indices as Hilbert coordinates; copying indices is a norm-one
coordinate restriction. It approximates the differentiated conditional
expectation itself, rather than differentiating an approximation to
its value. Thus, even when $u=u(X_0)$ is correlated with $(x,\pi)$,
\begin{equation}\label{eq:hidden-row-weighted-error}
 \norm{D_u(A-A_N)}_{\HS}^2
 =\sum_i|u_i|^2\norm{\operatorname{row}_i(A-A_N)}^2
 \leq C_F^2\varepsilon_N^2\norm{u}^2.
\end{equation}
This is the row-weighted truncation bound used in the covariance
calculation below.

Apply the operator estimate before opening auxiliary paths. Then
\Cref{lem:path-increment-approximation,prop:path-opened-derivative}
approximate the fixed mixed tensors by tensor networks in operator
norm and place their opened contractions in the path-matrix closure.
The retained matrices satisfy the entry and Hadamard rules of
\Cref{thm:path-calculus}. This concerns fixed capped histories after
opening; it gives no operator-norm approximation of an arbitrary
conditional expectation by finitely many Monte Carlo samples.

\begin{lemma}[A marked continuing derivative]\label{lem:hidden-marked-spine}
Fix an opened, capped derivative history of one smooth multitime
label. Suppose it has an entering derivative index, at least one
branch differentiating that label, and one derivative index continuing
to a later factor. Its temporal derivative tree has a unique path
between the entering and continuing indices. Absorb the part before
its first label attachment into the incoming propagation matrix, and
retain the part after its last attachment as a subsequent bounded
propagation. Let $M$ be the matrix arriving at the first attachment,
with label site as row and derivative index as column.

If $\Phi(M)$ is the matrix after the intervening history, then, in the
finite-tree comparison,
\begin{equation}\label{eq:hidden-marked-spine-bound}
 \norm{\Phi(M)}_{\HS}
 \leq C\norm{\diag(M)}+o(1)\norm{M}_{\HS}.
\end{equation}
Here $C$ depends on the fixed circuit and capped history. All memory
copies and all mixed differences of the label are included in this
history. The estimate does not replace a mixed label tensor by an
ordinary first Jacobian.
\end{lemma}

\begin{proof}
First fix a tensor-network approximation from
\Cref{lem:path-increment-approximation}. Before opening paths, use
the mixed-tensor budget to bound its operator errors, and remove
those errors in the prescribed successive approximation limit.
The temporal tree off the marked path consists of disjoint subtrees
ending at derivative leaves of the chosen label. Each attachment has
at least one such leaf. The corresponding maps are bounded tensor
products and compositions of capped path matrices and equality
copies; every equality copy has norm at most its fixed scalar label
bound. Pull their adjoints into the row tensor of the label circuit.
Repeated occurrences of a derivative index are handled by the
norm-one equality embedding. The resulting tensor $U_i$ on the $r$
attachment indices satisfies
\begin{equation}\label{eq:hidden-marked-row-tensor}
 \sup_i\sum_{a_1,\ldots,a_r}|U_i(a_1,\ldots,a_r)|^2\leq C^2.
\end{equation}
Use the operator norm of each flattened mixed tensor to bound each
of its rows in Euclidean norm. The off-path maps act on disjoint
groups of leaves, so their tensor-product operator norms give the
stated estimate without a dimension factor.

Suppress vertices on the marked path with no label attachment. The
matrices between its remaining consecutive attachments are bounded
path matrices $Q_1,\ldots,Q_{r-1}$. For $r\geq2$, put
\begin{equation}\label{eq:hidden-marked-kernel}
 L_i(a_1,a_r)=\sum_{a_2,\ldots,a_{r-1}}
 U_i(a_1,\ldots,a_r)\prod_{\ell=1}^{r-1}
                       (Q_\ell)_{a_\ell a_{\ell+1}}.
\end{equation}
Then $\Phi(M)_{i,b}=\sum_a M_{i,a}L_i(a,b)$. If $r=1$, this formula
means $L_i(a,b)=\1_{a=b}U_i(a)$.

The next step is to make every $Q_\ell$ in
\Cref{eq:hidden-marked-kernel} diagonal. Control one replacement at
a time: fix an edge $e$ and put $Q_e^{\rm off}$ there. With the
endpoints fixed, Cauchy--Schwarz bounds the squared kernel by
\[
 \left(\sum_{a_2,\ldots,a_{r-1}}|U_i(a_1,\ldots,a_r)|^2\right)
 \left(B_1\cdots B_e^{\rm off}\cdots B_{r-1}\right)_{a_1a_r},
\]
where $B_\ell=|Q_\ell|^{\circ2}$ entrywise. Every row sum and every
column sum of $B_\ell$ is at most $\norm{Q_\ell}_{\op}^2$.
Writing $\eta_e=\max_{a\ne b}|(Q_e)_{ab}|$, we have
$B_e^{\rm off}\leq\eta_e^2\1\1^t$ entrywise. The second factor is
therefore at most
\[
 \eta_e^2\prod_{\ell\ne e}\norm{Q_\ell}_{\op}^2.
\]
Sum over the endpoints and use \Cref{eq:hidden-marked-row-tensor}.
This gives
\[
 \sup_i\norm{L_i^{(e,\rm off)}}_{\HS}
 \leq C\eta_e\prod_{\ell\ne e}\norm{Q_\ell}_{\op}.
\]
The estimate remains valid after any other edges have been made
diagonal. Telescope over the finitely many edges and apply the entry
rule in \Cref{thm:path-calculus} to obtain
$\sup_i\norm{L_i-L_i^{\rm diag}}_{\op}=o(1)$. Its row action costs
at most $o(1)\norm{M}_{\HS}$.

After these replacements, diagonal marked-path edges force all
attachment indices to agree. The remaining action takes the form
$M\circ E$, where
\[
 E_{i,a}=U_i(a,\ldots,a)
                \prod_{\ell=1}^{r-1}(Q_\ell)_{aa}.
\]
This $E$ contains the entire label derivative subnetwork, including
its memory copies. It is the two-terminal network obtained by tying
all temporal attachment indices to the common source $a$, with the
label output $i$ as sink. Every derivative leaf is paired.
\Cref{lem:path-directed-network,prop:path-opened-derivative} place it
in the bounded path-matrix closure. The diagonal factors above may
first be replaced by their uniformly accurate diagonal predictions.
Consequently $\max_i|E_{ii}|\leq C$ and
$\max_{i\ne a}|E_{i,a}|=o(1)$. The estimate
\[
 \norm{M\circ E}_{\HS}^2
 \leq C^2\norm{\diag(M)}^2
       +\max_{i\ne a}|E_{i,a}|^2\norm{M}_{\HS}^2
\]
proves \Cref{eq:hidden-marked-spine-bound}. Here the entry estimates are used on fixed typical finite trees,
followed by their integrated comparison. Their scope does not include
uniformity over arbitrary states.
\end{proof}

\begin{figure}[htbp]
\centering
\begin{tikzpicture}[
  vertex/.style={circle,draw=black!65,fill=white,inner sep=2pt},
  edge/.style={-{Stealth[length=1.8mm]},draw=black!65},
  lab/.style={font=\small,fill=white,inner sep=1pt}]
 \node (in) at (0,0) {$M$};
 \node[vertex] (a) at (1.6,0) {$a_1$};
 \node[vertex] (b) at (4.7,0) {$a_2$};
 \node[vertex] (c) at (7.8,0) {$a_3$};
 \node (out) at (9.6,0) {$dq$};
 \draw[edge,line width=1pt] (in) -- (a);
 \draw[edge,line width=1pt] (a) -- node[lab,above] {$Q_1$} (b);
 \draw[edge,line width=1pt] (b) -- node[lab,above] {$Q_2$} (c);
 \draw[edge,line width=1pt] (c) -- node[lab,above] {suffix} (out);
 \node[draw=black!60,rounded corners=2pt,fill=black!3,
       minimum width=8.2cm,minimum height=8mm]
       (tensor) at (4.7,-1.4)
       {$U_i(a_1,a_2,a_3)$: pulled-back label tensor,
        $\|U_i\|_{\ell^2}\leq C$};
 \draw[edge] (a) -- (a |- tensor.north);
 \draw[edge] (b) -- (b |- tensor.north);
 \draw[edge] (c) -- (c |- tensor.north);
 \node[font=\small] at (4.7,-2.05)
      {nonempty label branches at every attachment};
\end{tikzpicture}
\caption{The derivative continuing into the later suffix defines
a marked path.  Side branches are absorbed into a label tensor
with bounded row square norm.  After branch-free segments have been
combined, every intervening $Q_r$ separates two nonempty sets of
label leaves.  Its off-diagonal contribution is small by
Cauchy--Schwarz and the row and column sums of the squared-entry
matrices.  Pinching all such segments ties the attachment indices
and returns the diagonal estimate for $M\circ E$.}
\label{fig:marked-continuation}
\end{figure}


\subsection{Proof of the averaging identity}
\label{subsec:hidden-averaging-proof}

\begin{proof}[Proof of \Cref{lem:hidden-averaging}]
Fix a summand of \Cref{eq:hidden-double-centered-error}, with
distinguished labels
$a=a_r-a_r^0$ and $b=a_q-a_q^1$, where $r<q$. By
\Cref{lem:hidden-small-trace}, approximate its path labels by fixed
bounded smooth multitime circuits, with every observed state inside
the label's time span. These circuits may contain mixed queries such
as $J[f(X_s)\odot X_r]$; no representation as a sum of separated
single-time products is required. The frozen-score construction of
\Cref{lem:path-likelihood-approximation} puts its smooth caps in the
class of \Cref{lem:path-mixed-tensors}; in particular, outside spins
multiplying an unbounded query are handled by its finite
sign-pattern decomposition. Before estimating derivatives, fix the circuit, its cap, and every
observation time. Keep both distinguished vertices exactly
conditionally centered for each approximant. Approximate the
nonspecial endpoint factors by fixed smooth one-state functions
before the derivative limits, and perform the field cutoff above.
Its threshold and the suffix approximants remain fixed throughout
the derivative comparisons.

For the exact right predictor, use a conditional path representation
instead of requiring a regular one-state recipe. Represent its value
by an independent ordinary stationary-kernel path from $X_t$, reading
the multitime evaluations in reverse order. Differentiate by propagating
the full mixed-slot recursion down this added causal branch.
Its tensor and path-network conclusions are those of
\Cref{lem:path-conditional-derivative,prop:path-opened-derivative}.
Keep the exact conditional subtraction before opening any of these
derivatives. The conditional estimates
\Cref{eq:hidden-conditional-op,eq:hidden-row-weighted-error} apply
there as well. The error in replacing a tracial predictor by this
exact prediction is handled before derivatives are opened, by
\Cref{lem:hidden-small-trace}.

Split the Wick string at $a$ and $b$ using
\Cref{eq:hidden-wick-split}. The subtractions involving $\tau(a)$ or
$\tau(b)$ vanish; errors from approximate trace centering are
controlled in the unsplit Wick form by
\Cref{lem:hidden-small-trace}. Absorb the outer strings into the
endpoint vectors. The left vector remains hidden by
\Cref{thm:path-calculus}; the right vector has bounded mean-square
unweighted Hessian by \Cref{eq:hidden-cutoff-hessian} and
\Cref{eq:hidden-matrix-product-rule}. It remains to prove
\begin{equation}\label{eq:hidden-covariance-target}
 \E u(X_0)^tD_a QD_b p(X_t)\longrightarrow0,
\end{equation}
where $Q$ is a positive-length Wick string with identity end labels.
The left centered label $a$ will give a martingale covariance. If its
paired derivative reaches $p$, the backward centering of $b$ will
cancel the term. If the derivative first hits a tail label, the
nonempty Wick prefix will supply the cancellation instead.

To prove the remaining bilinear limit, use conditional martingale
covariance up to the last evaluation in $a$. For scalar future
observables $F,G$, it reads
\begin{equation}\label{eq:hidden-martingale-covariance}
 \Cov_x(F,G)
 =2\E_x\int\sum_j w_j(X_s)
              d_jF_s(X_s)\,d_jG_s(X_s)\dd s,
\end{equation}
where $F_s,G_s$ are the conditional future expectations with their
already observed arguments held fixed. The integration stops once
$F$ is fully observed. To check the factor, a flip of $j$ has rate
$w_j/2$ and changes a scalar martingale by $-2x_jd_jF_s$; its product
bracket is $2w_jd_jF_s d_jG_s$. Since $a$ has exact forward conditional
mean zero, \Cref{eq:hidden-covariance-target} is the corresponding
covariance, with $u(X_0)$ fixed under the initial conditioning.

\Cref{eq:hidden-martingale-covariance} introduces a derivative
index in addition to the sites of the Wick string. Propagate the two
derivatives by the mixed-slot recursion and $dS_h=K_hd$, including
every exact product and memory-copy correction. First conditionally average the factor differentiating $a$ and
remove its truncation with \Cref{eq:hidden-row-weighted-error}.
On the other side, each terminal tensor term contains $p$ or $dp$
exactly once, because $p$ has only the terminal state slot. The bounded mixed label tensors and the Hilbert-space
amplification of \Cref{prop:path-rough} therefore propagate its
unweighted mean-square Hilbert--Schmidt budget. Tensor and
propagation errors are removed in these norms before paths are
opened. For the retained histories use
\Cref{prop:path-opened-derivative} to reduce their complete
derivative subnetworks to bounded path matrices. The two conditional
future expectations in
the covariance can then be represented by independent causal futures
from their joining state.

We now split the retained terms according to the derivative's first
destination: either it reaches $p$ before differentiating any tail
label, or it first differentiates such a label. The two cases use
different cancellations.

\emph{The derivative reaches $p$ without differentiating a tail
label.} Expand each terminal mixed tensor by the exact product rule,
and classify a term by whether any difference falls on a tail label.
For a label $c$ and a remaining factor $q$, we use
\[
 d_j(cq)=c\,d_jq+q\,d_jc-2x_j(d_jc)(d_jq).
\]
The term that bypasses $c$ retains its unshifted value; classify the
mixed correction with the first-hit terms. If no tail label is
differentiated, the only differentiated state slot is the terminal
slot of $p$. Any extra frozen-slot branch would differentiate a
label, as would a same-slot correction containing both a label
difference and $dp$; both belong to the second case. Thus the history
in the present case is unbranched and follows the common $K$
propagation through the entire observation grid. With the derivative as column index, its term is a
finite sum or integral of
\begin{equation}\label{eq:hidden-hessian-ending}
 \E\ip{dp(X_t)}{D_bQ^tD_u A}_{\HS},
\end{equation}
where $A$ is a bounded path matrix containing the derivative
propagation, the differentiated $a$-side, and the bounded rate
diagonal. Condition first on $X_t$, so the terminal Hessian is fixed when
we apply Cauchy--Schwarz. We must estimate the conditional matrix
\[
 H(X_t)=\E[D_bQ^tD_u A\mid X_t].
\]
Take two conditionally independent copies of the entire path system
given $X_t$. Directly expanding its Hilbert--Schmidt square gives
\begin{equation}\label{eq:hidden-hadamard-pair}
 \E\norm{H(X_t)}_{\HS}^2
 =\E\,u_1^t\left[
       (Q_1D_{b_1b_2}Q_2^t)\circ(A_1A_2^t)
                    \right]u_2.
\end{equation}
Reversibility and the tree comparison identify this system as a causal
tree rooted at $X_t$. Each conditional copy shares that one root;
we do not introduce a new bridge conditioned on two endpoints.

Apply the Hadamard rule in \Cref{thm:path-calculus}. Operator-small
errors and visible dyads have zero bilinear limit between these
hidden vectors. The remaining diagonal is the product of the two
diagonal predictions. The prediction of
$Q_1D_{b_1b_2}Q_2^t$ is zero. Indeed, the Wick strings admit no internal
return, so the two equal-length strings must pair fully across each
other. Their innermost contraction is $\tau(b_1b_2)$. Under the
stationary kernel, the two copies are conditionally independent and
$\E[b_{1,i}b_{2,i}\mid X_t]=0$ at every site. The zero limiting pair
trace transfers to the actual, possibly tilted, reversed tree by
\Cref{thm:path-calculus}; exact finite-dimensional stationary
centering is not asserted for that actual reversed kernel. Thus the
right-hand side of \Cref{eq:hidden-hadamard-pair} tends to zero, and
the Hessian budget proves that \Cref{eq:hidden-hessian-ending}
vanishes.

This cancellation requires common subdivisions and identical
truncated propagation kernels for the actual and predictor terms.
Keep every evaluation time in the subdivision, even when the factor
evaluated there is absent from a particular term. In the present case no derivative falls
on $b$, so its actual and predictor pieces have the identical matrix
$A$. Recombine them into the exactly centered $b$ before using
\Cref{eq:hidden-hadamard-pair}. Independent predictor simulation
branches may then be used in the two copies. Separate truncations
followed by separate Hadamard estimates would not establish the
vanishing of $\tau(b_1b_2)$.

\emph{A tail label is differentiated.} Stop at the first such vertex
from the left, including a product-rule term that differentiates the
label while also continuing to the right. If an earlier contraction
has made this vertex scalar, its coefficient is a deterministic
limiting trace from the word calculus, not a state-dependent empirical
trace, and is not differentiated. The remaining prefix retains its
Wick inclusion--exclusion.

The stopping rule must be fixed before opening derivatives of the
later suffix. The surviving actual labels have ordered spans with
disjoint interiors: an adjacent return removes the interior excursion
and merges the two neighboring labels, whose convex span contains no
surviving vertex span in its interior. Repeating returns preserves
this property.
The endpoint predictions in the double-centered string occur only
in its terminal suffix, apart from the initial predictions already
absorbed into $u$. If the first differentiated actual label $c_i$ has
span ending at $v$, every later actual label is observed after $v$,
and every later endpoint predictor is a function of $X_t$. For each
fixed spatial suffix, form its exact conditional value
\begin{equation}\label{eq:hidden-stopped-suffix}
 q_i(z)=\E[\text{later spatial suffix ending in }p\mid X_v=z].
\end{equation}
The Markov property makes the local observable exactly
$c_i(\text{path up to }v)q_i(X_v)$. Spatial matrix contractions in
that suffix are included in $q$; they need not be diagonal.

Before reopening the suffix, we need its $L^2$ value and mean-square
Hilbert--Schmidt derivative to be bounded. Prove both by working
backwards from the bounds for $p,dp$, using the bounded mixed label
tensors and exact product and copy rules at each ordered span. For the empty
derivative-index set, this means multiplication by the bounded label
$c$, with bound $C_{\varnothing}=\norm c_\infty$. The single
terminal state slot contributes the value or the first difference of
the current suffix once. Their tensor Hilbert--Schmidt norms are
bounded by the row tensor bounds of the labels times the corresponding
suffix norm. The Hilbert-space amplification of $S,K$ and their
common tails propagates these bounds to the preceding span, while a
spatial matrix costs its bounded operator norm. For a terminal
predictor, the ordinary one-state product rule uses its bounded
conditional derivative operator. This finite backward induction uses no second difference of $p$ in
its original state slot and does not require $q$ to be an exact
gradient.

Now apply the local temporal recursion to $c_iq_i$, retaining $q$
and $dq$ as unopened terminal objects. A term that differentiates $c$ either terminates in
its label tensor, leaving $q$ as a value, or has exactly one marked
leaf carrying $dq$. Every other derivative leaf belongs to this
single label $c$, including all its mixed-slot and memory-copy
leaves. Thus the latter term has precisely the marked history in
\Cref{lem:hidden-marked-spine}, illustrated in
Figure~\ref{fig:marked-continuation}. The later differentiated labels remain inside $dq$ until after this
local estimate. Conditioning in this order keeps all side leaves at
the same spatial output site, which a simultaneous expansion of the
later labels would not ensure.

For a first hit on a terminal predictor $g(X_t)$, apply the exact
one-state product rule to $D_gq$ before opening the auxiliary causal
branch representing $dg$. Its terminating term uses the whole
two-terminal matrix $dg$; its continuing correction is the
one-attachment case of \Cref{lem:hidden-marked-spine}. Actual and
predictor pieces of $b$ can be treated separately in this first-hit
case, since the cancellation below uses the Wick prefix, not a
later centering. This completes the stopping rule for every possible
first differentiated vertex.

Up to bounded row multipliers at that vertex, and with an unbranched
temporal segment before its first label attachment absorbed into
$A$, the incoming matrix is
\begin{equation}\label{eq:hidden-first-vertex-matrix}
 M=P^tD_uA,
\end{equation}
where $P$ is a positive-length Wick prefix ending in identity. Its
Hilbert--Schmidt norm is bounded by
$\norm{P}_{\op}\norm{A}_{\op}\norm{u}$, whereas its diagonal satisfies
the exact identity
\begin{equation}\label{eq:hidden-diagonal-cancellation}
 \diag(M)=(P\circ A)^tu.
\end{equation}
The diagonal prediction of $P$ is zero, since it is a nonempty Wick
string. The Hadamard rule and hiddenness hence show that
$\E\norm{\diag(M)}^2\to0$.

For a continuing term, use \Cref{lem:hidden-marked-spine} on the
whole derivative subnetwork of $c$, including more than its first
Jacobian. This gives
\begin{equation}\label{eq:hidden-schur-small}
 \E\norm{\Phi(M)}_{\HS}^2
 \leq C\E\norm{\diag(M)}^2
       +o(1)\E\norm{M}_{\HS}^2
 \longrightarrow0
\end{equation}
in the typical weighted comparison. The exact $-2x_j$ factors in
the product and copy rules are bounded labels on the equality
forks of this history. An unbranched segment after the last label
attachment is a subsequent bounded propagation. Pairing with $dq$
at the stopping boundary now gives zero by Cauchy--Schwarz and its
Hilbert--Schmidt budget. Equivalently, later suffix operations may
then be opened; their bounded Hilbert--Schmidt maps preserve this
zero estimate.

If the derivative terminates at its first differentiated vertex,
let $E$ be the matrix of the complete terminating label subnetwork.
All its mixed leaves and memory copies have been contracted.
By \Cref{prop:path-opened-derivative} this is a bounded path matrix.
Contract its incoming index against $M$. In indices, the vector left
at that vertex is
\[
 \sum_j M_{ij}E_{ij}
 =\sum_h P_{hi}u_h(AE^t)_{hi}
 =\bigl[(P\circ(AE^t))^tu\bigr]_i.
\]
It tends to zero by the same Hadamard rule, again because the
prediction of $P$ is zero. Bounded row or column labels can be included
in this identity without changing its conclusion. Pair with the
bounded $L^2$ value of $q$ from \Cref{eq:hidden-stopped-suffix}.
The terminating and continuing cases together cover every term and
correction supplied by the exact product and copy rules.

The finitely many deterministic observation times may be separated
within their prescribed spans before taking limits. Integration times
coinciding with them have zero Lebesgue measure. The ordinary rough
bounds control the limits and all discarded propagation tails. We
have proved \Cref{eq:hidden-covariance-target}; summing
\Cref{eq:hidden-double-centered-error} proves
\Cref{eq:hidden-averaging}.
\end{proof}

\subsection{Likelihood caps and endpoint norm estimates}
\label{subsec:hidden-caps}

Averaging is now available. Apply \Cref{lem:hidden-averaging} to the
path expansion of $K_t$. By \Cref{prop:path-rough}, deleting hops of
order $k>A_Lt$ and elementary likelihoods above $\exp(C_Lt)$ costs
$O(e^{-Lt})$ in the endpoint $L^2$ comparison for any prescribed $L$.
Hidden orthogonality removes the visible ranks. We are left with
partially paired Wick strings having $k=O(t)$, for which the main
issue is the supremum of an endpoint predictor at a Fock cut.

At each trunk vertex, the label is a nonnegative product of row
likelihoods on a deterministic union of waiting intervals and factors
between zero and one. Every vertex after the first also carries its
arrival factor $v$. The likelihood product is therefore a single
deterministically masked likelihood. For a span of length $d$, its raw
trace moments have the rough fixed-order bound $e^{C_pd}$, whereas
the projected row moments satisfy
\begin{equation}\label{eq:hidden-projected-label-moments}
 \max_i\norm{\E_x[a_i\mid X_d]}_{L^p(X_d)}
 \leq e^{o_p(d)}
\end{equation}
at typical starts, by \Cref{thm:row-projected-moments}. Restricted
labels are handled by domination by the unrestricted likelihood.
The forward predictor is at most one. Under stationarity the reverse
predictor has spatial moments of every fixed order bounded by
$e^{o_p(d)}$. These statements concern deterministic masks, with
uniform constants; the quenched exceptional mass is integrated over
simplex times as in \Cref{thm:row-projected-moments}. A
disorder-selected worst mask is not needed.

The actual center label left by \Cref{eq:hidden-averaging} can be
estimated directly by a projected row moment. If $U=U(X_0)$ and $V=V(X_t)$ are the vectors
after the two endpoint strings have been applied, set
$g_i(X_t)=\E_x[a_i\mid X_t]$. Row Cauchy--Schwarz gives
\begin{align}
 \abs{\E_x\sum_iU_i a_iV_i}
 &\leq\sum_i |U_i|\norm{g_i}_{L^2(X_t)}
                         \norm{V_i}_{L^2(X_t)}\notag\\
 &\leq\bigl(\max_i\norm{g_i}_{L^2(X_t)}\bigr)
               \norm{U}\bigl(\E_x\norm{V}^2\bigr)^{1/2}.
 \label{eq:hidden-center-row-cs}
\end{align}
Average in the initial state and use Cauchy--Schwarz again. The
polynomial endpoint bounds remove the exponentially atypical starts.
Thus the center costs $e^{o(t)}\norm{U}_2\norm{V}_2$, and an
adjacent-pair subtraction at the center costs only $|\tau(a)|$.

The remaining task is to bound the endpoint predictor suprema in
\Cref{eq:hidden-cut-bound} at a cost of only $e^{\epsilon t}$.
An untruncated raw likelihood has no operator bound available here.
We first fix a raw likelihood cutoff using one fixed rough moment.
Only afterwards do we choose the higher projected moment needed to
make the predictor supremum loss small. Keeping these choices separate
prevents the high-order rough constants from entering that loss. Begin by fixing the error rate $L$ and the rough
path truncations. Restrict the composite labels to the box
$\max_j\log^+ a_j\leq H_0$, with $H_0=C't$ and $C'$ to be chosen.
Outside that box use the dyadic annuli
$H<\max_j\log^+a_j\leq2H$, $H=2^mH_0$.
Interpret this maximum as a tensor box in the separate label
arguments. Multilinearity decomposes it into a separate restriction
at each label; it is not a pointwise maximum at one common site.
Telescope the difference between two boxes by the first vertex whose
label exceeds the lower threshold. This uses at most $k+1$ terms
per annulus. The elementary caps leave only $O(\log(t+2))$ annuli
because the number of intervals is $O(t)$.

In one such box term, continue to write $a_j$ for its restricted
label, so $0\leq a_j\leq e^{2H}$. Write $h_j^s=a_j^s$, $s=0,1$,
for its exact stationary endpoint predictors. For a parameter
$0<\epsilon'\leq1$, set
\[
 \widetilde h_j^s=\min\{h_j^s,e^{\epsilon'H}\},\qquad
 e_j^s=h_j^s-\widetilde h_j^s,
 \qquad c_j^s=a_j-h_j^s.
\]
At the designated exceptional vertex set both $\widetilde h_j^s=0$
instead. Denote by $\delta_{\rm raw}$ a bound for the trace two-norm
of its discarded predictors, and by $\delta_{\rm pred}$ a bound for
the other clipping errors. In the main box there is no designated
vertex and $\delta_{\rm raw}=0$.

We next realize the exact capped endpoint functions in the
finite-tree calculus. This step must avoid using a small-trace
replacement at an arbitrary exterior position. Fix $M=e^{2H}$ and an endpoint state
$y$. For $h(y)=\E[a\mid y]$, take $N$ conditionally independent
stationary-kernel branches with labels $a^{(1)},\ldots,a^{(N)}$ in
$[0,M]$, and put
\[
 Z_N=\frac1N\sum_{r=1}^N a^{(r)},\qquad
 g_N(y)=\E[f(Z_N)\mid y],
\]
where $f(z)=\min\{z,e^{\epsilon'H}\}$ or
$f(z)=(z-e^{\epsilon'H})_+$. Both maps are $1$-Lipschitz, so
conditional variance and Jensen give
\begin{equation}\label{eq:hidden-cap-branch-approximation}
 \begin{aligned}
 \sup_y\max_i|g_{N,i}(y)-f(h_i(y))|&\leq M/\sqrt N,\\
 \sup_y\frac1n\sum_i
   \E[|f(Z_{N,i})-f(h_i(y))|^2\mid y]&\leq M^2/N.
 \end{aligned}
\end{equation}
For ordinary prediction, take $f(z)=z$. Assign an independent branch
group to each occurrence of a conditional average in a word.
Multilinearity represents the resulting products and trace
contractions as expectations on a finite causal tree. The first bound
controls the diagonal error uniformly in operator norm, so it can be
removed from any bounded fixed word. The second recovers the trace
norms from the sampled labels. Both the cap labels and the error
labels retain their respective supremum bounds.
If needed, the scalar cap is uniformly smoothed on $[0,M]$ before
using the tree calculus. These limits are taken at fixed $t,H$
before the large-time estimate. They require neither a derivative
bound for the exact endpoint predictor nor a high raw likelihood
moment. True centered factors retain their exact subtractions.

Using the modified predictors in \Cref{eq:hidden-averaging} has the
exact residual
\begin{equation}\label{eq:hidden-modified-residual}
 \sum_{i<j}B_k\bigl(\widetilde h^0_{<i},\ c_i^0+e_i^0,
            \ a_{i<q<j},\ c_j^1+e_j^1,\ \widetilde h^1_{>j}\bigr).
\end{equation}
Expand the residual into three cases. With two true centered
factors, the averaging proof gives zero; that proof permits bounded
endpoint-only factors outside the distinguished vertices. With one
error and one centered factor, every Fock cut away from the error
pays its trace two-norm. At the error itself, the centered factor
lies on its appropriate pure side, and the single-star argument of
\Cref{lem:hidden-small-trace} gives zero. With two errors, at least
one is off the cut and pays its trace two-norm. If the other is at
the cut, retain its supremum bound $e^{2H}$, including its possible
action on the rare base.

All remaining trace factors are bounded using only the raw second
moment. Their product costs $e^{Ct}$: the waiting masks across the
vertices use disjoint waiting intervals. The fixed powers per
vertex, pairing choices, cuts, and telescoping terms cost at most
another $e^{Ct}$ because $k\leq A_Lt$. Consequently, after the
dimension and approximation limits, the total modified-averaging
error is bounded by
\begin{equation}\label{eq:hidden-cap-error}
 e^{C_0t+2H}
       (\delta_{\rm raw}+\delta_{\rm pred})
       \norm{u}_2\norm{p}_2.
\end{equation}
Only the rough truncations and second moments enter the fixed
constant $C_0$. It is independent of $\epsilon'$ and of the high
predictor moment order to be chosen below. The successive limits
have already removed the derivative-approximation constants from
the exact centered cancellation, so they do not enter $C_0$.

Treat the designated raw tail first. For each fixed $r>0$, the
uniform rough likelihood moment yields
\begin{equation}\label{eq:hidden-raw-tail}
 \max_i\E_x[a_i^2\1_{a_i>e^H}]
 \leq e^{-rH}\max_i\E_x a_i^{r+2}
 \leq\exp(C_{r+2}t-rH).
\end{equation}
The bound holds at every starting state. Conditional Jensen therefore
gives both
\begin{align}
 \delta_{\rm raw}
 &\leq\exp\left(\tfrac12C_{r+2}t-\tfrac r2H\right),\notag\\
 \max_i\norm{\E_x[a_i\1_{a_i>e^H}\mid X_t]}_{L^2(X_t)}
 &\leq\exp\left(\tfrac12C_{r+2}t-\tfrac r2H\right).
 \label{eq:hidden-exceptional-center}
\end{align}
Stationarity gives the first estimate, which also bounds the scalar
adjacent-pair trace by Cauchy--Schwarz. The second supplies the
uniform row bound at the exceptional actual center. Both use the
same fixed raw moment; no high projected moment is needed here.

In an annulus, setting both designated predictors to zero kills every
term in the modified averaging formula except the actual center at
that vertex. Split there by \Cref{eq:hidden-wick-split}. Each of its
two endpoint strings pays at most one supremum
$e^{\epsilon'H}$ by \Cref{eq:hidden-cut-bound}. Apply
\Cref{eq:hidden-center-row-cs,eq:hidden-exceptional-center} to the
center. Its contribution, its adjacent-pair subtraction, and the raw
part of \Cref{eq:hidden-cap-error} are all bounded by
\begin{equation}\label{eq:hidden-raw-annulus-cost}
 \exp\left[\left(C_0+\tfrac12C_{r+2}\right)t
                      -\left(\tfrac r2-2\right)H\right]
       \norm{u}_2\norm{p}_2,
\end{equation}
where we used $\epsilon'\leq1$. Fix, for example, $r=12$, and then
choose $C'\geq1$ so that
\begin{equation}\label{eq:hidden-raw-cap-choice}
 4C'\geq C_0+\tfrac12C_{14}+L+3.
\end{equation}
For all $H\geq C't$, \Cref{eq:hidden-raw-annulus-cost} is at most
$e^{-(L+3)t-4(H-C't)}\norm{u}_2\norm{p}_2$. This fixes $C'$ without
reference to the desired small supremum loss.

With the raw cutoff fixed, prescribe $\epsilon>0$ and set
$\epsilon'=\min\{1,\epsilon/(2C')\}$. The two endpoint suprema in
the main box now cost at most $e^{\epsilon t}$. Choose the high
predictor-tail order $m$ only after making this choice. For every ordinary
predictor clipping error, domination by the unrestricted predictor
and the projected moment theorem give
\begin{equation}\label{eq:hidden-predictor-tail}
 \begin{aligned}
 \delta_{\rm pred}^2
 &\leq\max_{j,s}\tau((h_j^s)^2\1_{h_j^s>e^{\epsilon'H}})\\
 &\leq e^{-m\epsilon'H}\max_{j,s}\tau((h_j^s)^{m+2})
 \leq\exp[-m\epsilon'H+o_{m+2}(t)].
 \end{aligned}
\end{equation}
Here the mask can be viewed on the full interval $[0,t]$; a
restricted label has predictor bounded by that of the unrestricted
likelihood. Choose $m$ sufficiently large that
\begin{equation}\label{eq:hidden-predictor-cap-choice}
 \frac{m\epsilon'}2-2\geq\frac{C_0+L+4}{C'}.
\end{equation}
The predictor part of \Cref{eq:hidden-cap-error} is then at most
$e^{-(L+3)t}$ for all sufficiently large fixed $t$, with a further
decaying exponential in $H-C't$. This holds in the main box and in
every annulus. Summing the annuli therefore costs $O(e^{-Lt})$.

The raw constant $C_{m+2}$ is absent from
\Cref{eq:hidden-predictor-tail,eq:hidden-predictor-cap-choice} because
the projected limiting moment theorem has already removed the rough
moments used for fixed-time uniform integrability. Increasing the
fixed order $m$ may increase the eventual threshold in $t$, while
leaving $C'$ unchanged. Only the one or two distinguished errors in
each residual term use this high order; every other trace factor
still uses order two.

In the surviving main-box terms the actual center costs $e^{o(t)}$
by \Cref{eq:hidden-center-row-cs}. The endpoint-only subtraction
terms in \Cref{eq:hidden-averaging} satisfy the direct cut bound with
one predictor supremum. Thus the combined loss is
$e^{\epsilon t+o(t)}$, and only a fixed number of trunk vertices
lose their trace two-norm estimate. The order of choices is the
fixed raw order, then $C'$, then $\epsilon'$, then the high projected
order $m$.

Apply these endpoint polynomial norm bounds after removing visible
components. If a visible component is created on the right, it
remains orthogonal to the hidden left vector after the remaining
path operations. Prove this first for a finite endpoint recipe
projection, then increase the basis using the tree Gram comparison.
Thus the bound used here is the hidden word domination of
\Cref{thm:path-calculus}; no unrestricted operator comparison is
applied to a visible component.

\subsection{Trunk renewal, exceptional vertices, and temporal boundaries}
\label{subsec:hidden-renewal}

The remaining task is to sum the surviving strings. Fix
$z<\lambda_{\rm sp}$ with room to increase the parameter slightly
without reaching the edge. Every removed side excursion has a
positive, completely paired stationary weight. Recall the rooted
paired kernel $R$ from \Cref{prop:spin-spectrum}, and put
\[
 \widehat R(z)=\int_0^\infty e^{zd}R(d)\dd d.
\]
The two inequalities in \Cref{eq:spin-conditions} are
\begin{equation}\label{eq:hidden-renewal-conditions}
 \beta\widehat R(z)<1,
 \qquad 1-z-\beta^2\widehat R(z)>0.
\end{equation}
The first inequality makes interior runs summable. The second
controls vertices whose arrival or trace weight has been lost.

An ordinary trunk vertex contributes the trace two-norm of its
initial or terminal predictor. For a fixed shape with bounded span
$d$, this norm approaches the traced weight $\tau(a)$ as the vertex
moves away from the corresponding global time boundary. To justify
that replacement, use the following equilibrium comparison. If $F_i$ is
the conditional expectation of the vertex functional given the state
at the beginning of its span, the gap yields
\[
 \frac1n\sum_i\norm{S_sF_i-\mu F_i}_{L^2(\mu)}^2
 \leq e^{-2\gamma s}\frac1n\sum_i\norm{F_i-\mu F_i}_{L^2(\mu)}^2.
\]
The two independent stationary-sample comparison of
\Cref{thm:path-calculus} gives
\[
 \frac1n\sum_i(\mu F_i)^2\longrightarrow\tau(a)^2.
\]
Taking square roots bounds the predictor trace norm by
$\tau(a)+o_s(1)$; the reversed comparison is identical. Fixed-shape
moment bounds justify the displayed averages. Box restrictions can
be dropped by positivity. For long spans, the reverse predictor has
cost $e^{o(d)}$ by \Cref{thm:row-projected-moments}, independently of
its location within the global interval, and the forward predictor
has cost at most one.

For exceptional vertices, discard the arrival factor $v$ and charge
one for the remaining label. The resulting positive rooted kernel
admits the following majorant. Let $B(d)=e^{-d}\1_{d\geq0}$, and let
$E(d)=\beta^2R(d)$. Arbitrarily many excursions, separated by bare
waits, give
\begin{equation}\label{eq:hidden-exceptional-kernel}
 V=B*\sum_{m\geq0}(E*B)^{*m},
 \qquad
 \widehat V(z)
 =\frac{1}{1-z-\beta^2\widehat R(z)}<\infty.
\end{equation}
The zeroth convolution power is the unit mass at zero. Thus $V$ is a
majorant for a vertex missing its arrival weight, and for the
exceptional vertices whose trace estimates were lost above. A
subexponential span factor is harmless: for every sufficiently small
$\delta>0$ it is bounded by $C_\delta e^{\delta d}$, and
$\widehat V(z+\delta)<\infty$. In particular the corresponding
$e^{zd}$-weighted tails tend to zero.

Return to ordinary vertices. Their arrival weight and stationary
trace give exactly $R$, with an additional factor $\beta$ for the
next unpaired edge. First restrict to bounded spans, finitely many
excursions, and bounded individual excursion lengths. For these
fixed shapes, the gap estimate is uniform beyond boundary distance
$T$, with error tending to zero as $T\to\infty$. The exceptional
majorant, using a small amount of slack, controls the discarded
shapes. Monotone truncation followed by dominated
convergence therefore provides an interior-vertex majorant $G_T$
such that
\begin{equation}\label{eq:hidden-interior-mass}
 \int_0^\infty e^{zd}G_T(d)\dd d
 \leq\beta\widehat R(z)+o_T(1)<1
\end{equation}
when $T$ is sufficiently large. Equivalently, first fix a shape and take the supremum of its predictor
trace bound over boundary distances at least $T$; then apply the
same truncation and domination. This order does not require uniform
approximation over an unbounded family of shapes.

Vertices lying within distance $T$ of time $0$ or $t$ use the crude
bound. The total sum of any number of hops contained in such a
bounded interval is finite: the order-$j$ contribution has the
simplex factor $T^j/j!$, and the fixed-time rough bounds give a sum
bounded by $e^{C_T}$. There are at most two additional vertices whose
spans cross the edges of these boundary intervals, since the trunk
spans have disjoint interiors. Those use $V$. The initial vertex
without an arrival factor, the chosen center, and the fixed number
of vertices at which a supremum was used also use $V$, with the
global multiplier $e^{O(\epsilon)t+o(t)}$ already identified.

We can now sum an entire term as an ordered concatenation of two
bounded-time boundary runs, a bounded number of exceptional kernels,
and interior runs. By \Cref{eq:hidden-interior-mass}, the weighted
mass of an interior run is a geometric sum. Choosing the center,
the cuts, and the finitely many exceptional locations contributes
only a polynomial in the number $j$ of interior steps. For a number
$q<1$ supplied by \Cref{eq:hidden-interior-mass},
\[
 \sum_{j\geq0}(j+1)^Cq^j<\infty.
\]
Consequently the entire positive majorant has finite Laplace mass at
$z$, up to $e^{O(\epsilon)t+o(t)}$.

To finish the renewal estimate, we must pass from finite Laplace
mass to a pointwise-in-time bound. For this, distinguish a boundary
or exceptional root that has a bare initial wait. In \Cref{eq:hidden-exceptional-kernel}
write $V=B*\nu$, where
$\nu=\sum_{m\geq0}(E*B)^{*m}$ has finite $z$-weighted mass. Since
$z<1$,
\begin{equation}\label{eq:hidden-density-bound}
 \sup_{d\geq0}e^{zd}V(d)
 \leq\sup_{r\geq0}e^{zr}B(r)
                    \int e^{zu}\nu(\dd u)<\infty.
\end{equation}
Subexponential factors are again paid using a slightly larger
parameter. If a boundary restriction cuts this root span short,
keep its bare first wait and sum all remaining vertices within that
bounded boundary interval by the same simplex estimate. Thus one
factor in each concatenation has bounded weighted density, while
all other factors have finite weighted mass. The convolution
inequality
\[
 \norm{f*g}_{\infty,z}
 \leq\norm{f}_{\infty,z}\norm{g}_{1,z},
 \quad
 \norm{f}_{\infty,z}=\sup_{d\geq0}e^{zd}|f(d)|,
 \quad
 \norm{g}_{1,z}=\int e^{zd}|g|(\dd d),
\]
gives the pointwise bound $C_z e^{-zt}$ for the summed majorant.

\begin{proof}[Completion of the proof of
\Cref{thm:hidden-propagation}]
Normalize the limiting endpoint second moments to one, using the
rough comparison if an endpoint norm vanishes. For each fixed $t$,
first truncate the hops and elementary likelihoods. Then apply
\Cref{lem:hidden-averaging}, remove the likelihood annuli, and sum
with the positive renewal estimate just proved. It gives, for every
$z'<\lambda_{\rm sp}$,
\[
 \limsup_{n\to\infty}|\E u^tK_t p|
 \leq C_{z',\epsilon}
       e^{-z't+O(\epsilon)t+o(t)}+O(e^{-Ht}).
\]
The derivative and field-cutoff errors were removed before the
large-time limit. Their constants therefore impose no growth
condition on \Cref{eq:hidden-energy-hypothesis}. Choose $z'>z$, then the predictor
cap parameter $\epsilon$ sufficiently small and the tail rate $H$
sufficiently large. The result is bounded by $C_ze^{-zt}$, enlarging
$C_z$ on a bounded interval of times using
\Cref{prop:path-rough}. Restoring the endpoint norms proves
\Cref{eq:hidden-propagation}.

Finally, the exceptional starting-state sets in the row estimates
have exponentially small integrated mass for the deterministic
simplex masks. Fubini and the polynomial endpoint bounds extend
the proof to the integrated comparisons in the statement. Every
replicated diagram used above is a causal tree; no step calls for a
uniform estimate on an adaptively selected bridge or mask.
\end{proof}
